\documentclass[12pt]{article}
\usepackage[utf8]{inputenc}
\usepackage[T1]{fontenc}
\usepackage{hyperref}
\title{Tagbu Language Resources for K-12 Curriculum Development}
\author{Compiled by Claude Sonnet 5 (claude-sonnet-5) for Andr\'as Kornai}
\date{}
\begin{document}
\maketitle

\section{Introduction}
Tagbu (also spelled Tagbo or Tagba; ISO 639-3 tbm; Glottocode tagb1255) is a Ubangian language of the northeastern Democratic Republic of the Congo. Wikipedia reports about 17,000 native speakers (2002 estimate). The sources found do not clearly show that the language is extinct. Very little material exists in the search results, and no educational resources of any kind were found.

\section{Primary Educational Resources}
None found.

\section{Linguistic and Reference Materials}
\begin{itemize}
\item Wikipedia, Tagbo language: classification and speaker estimate. \url{https://en.wikipedia.org/wiki/Tagbo_language}
\item Glottolog, Tagbu languoid record. \url{https://glottolog.org/resource/languoid/id/tagb1255.bigmap.html}
\item Numeral systems of Channel Islands/Africa database entry for Tagbu (Max Planck Institute for Evolutionary Anthropology). \url{https://lingweb.eva.mpg.de/channumerals/Tagbu.htm}
\end{itemize}

\section{Religious and Community Resources}
Joshua Project notes that Bible resources in the language are limited and lists the Tagbo people in the DRC. No scripture was identified.
\begin{itemize}
\item Joshua Project, Tagbo in the DRC. \url{https://joshuaproject.net/people_groups/15176/CG}
\end{itemize}

\section{Digital and Audio Resources}
None found.

\section{Curriculum Development Context}
No orthography, primer, dictionary or school programme was found. A heritage or revitalization curriculum would have to start from fieldwork or from archival material not located in this search. The numeral list and the Glottolog references are only starting points for finding descriptive sources.
\end{document}